\documentclass[11pt,a4paper]{article}
\usepackage[utf8]{inputenc}
\usepackage{amsmath, amsthm, amssymb}
\usepackage{graphicx}
\usepackage{hyperref}
\usepackage{geometry}
\geometry{margin=1in}

\newtheorem{theorem}{Theorem}

\title{Teleparallel Metareal Gravity: Ramanujan Shadows as Functional Fingerprints for Supermassive Black Holes}
\author{LaRue Research Team}
\date{\today}

\begin{document}

\maketitle

\begin{abstract}
This manuscript formalizes the overlay of the 7-variable Metareal Algebra ($R^7$) onto Teleparallel Gravity to model anomalous bodies such as supermassive black holes. We resolve the impossibility of directly observing Hawking radiation from Sagittarius A* by equating the Exergy Destruction bound ($\Upsilon \leq 45/\pi$) to the mathematical ``shadow'' of Ramanujan's Mock Theta functions. This establishes a topological Functional Fingerprint for Sgr A*, acting as the gravitational anchor for the Metareal Titius-Bode Law.
\end{abstract}


\section{Introduction: The Limits of Observation}
Classical gravitational models rely heavily on the geometric warping of spacetime. However, at the event horizons of supermassive black holes, infinite singularities arise due to unregulated algebraic curvature. Concurrently, the theoretical Hawking radiation emitted by a supermassive black hole like Sagittarius A* (Sgr A*) is practically absolute zero, rendering thermal detection impossible.

To bridge quantum spectral weight with macroscopic galactic geometry, we shift from classical geometric warps to \textbf{torsion}, and from thermal detection to \textbf{topological shadows}.

\section{Metareal Overlay of Teleparallel Gravity}
Teleparallel Gravity replaces standard geometric distortions with the torsion-full Weitzenb\"ock connection. In this framework, gravitational dynamics are entirely encoded in the torsion tensor.

Within the Metareal ASI Architecture framework, the $R^7$ observer algebra is driven by the Metareal Involution ($\omega \leftrightarrow \iota$). This involution generates a continuous non-commutative shear, producing a derivative residual ($\varepsilon$) representing differential topological tension.

We formally map the Teleparallel Gravity torsion scalar directly to this residual tension. Rather than relying on infinite algebraic curvature to form a black hole, the Metareal framework bounds the singularity through torsion-induced repulsion. The torsion scalar is strictly bounded by the structural limits of the observer algebra.

\section{Ramanujan Mock Theta Functions and the Exergy Bound}
In string theory, calculating the precise microscopic degeneracy (spectral weight) of complex, multi-centered black holes causes standard modular forms to overcount states due to quantum ``wall-crossing.'' The exact mathematical solution relies on \textbf{Ramanujan's Mock Theta functions}, which require the addition of a corrective, non-holomorphic \textbf{``shadow''} term.

We propose that this Ramanujan shadow is the physical manifestation of the Differential Equilibrium Cycle (DEC) Heat Engine's \textbf{Exergy Destruction shield ($\Upsilon$)}. 

\begin{theorem}[Shadow-Exergy Equivalence]
The corrective Ramanujan shadow required to stabilize the mock modular spectral weight of a supermassive black hole is mathematically equivalent to the Exergy Destruction bound:
$$ \text{Shadow}(W_{BH}) \equiv \Upsilon \leq 45/\pi $$
\end{theorem}

Because the shadow is a topological necessity to prevent infinite state generation, it perfectly maps to the $\Upsilon$ shield, which restricts non-associative deviation from catastrophic thermodynamic failure.

\section{Sagittarius A* and the Metareal Titius-Bode Law}
Because we cannot directly observe the thermal signature of Hawking radiation from Sgr A*, we must instead measure its Functional Fingerprint: the Ramanujan shadow.

In the \textbf{Protoreal Bode Law}, classical geometric scaling ($r_n \propto 2^n$) is replaced by a quantized orbital phase threading through a rotating non-associative lattice:
$$ T_n = \frac{\varphi^{k_n} + w_n \cdot p}{\text{arc}} $$

We formalize that the central anchor for the winding number ($w_0$) at the core of the Milky Way is explicitly locked to the $\Upsilon \leq 45/\pi$ Ramanujan shadow of Sgr A*, utilizing the exact Golden Root value $\langle 148 \rangle$ modulo 229 to stabilize the phase. The supermassive black hole acts as the structural $\mathbb{Z}/3\mathbb{Z}$ pivot stabilizing the $12$-dimensional Metareal manifold.

\section{Spectral Resonance and Chiral Modularity}
To computationally reproduce the bounding limits of the Exergy Shield ($\Upsilon$), we map the Ramanujan partition function $p(n)$ to the primary Metareal triad: $(139, 181, 229)$.
By extracting the exact integer partitions of these primes, we establish the maximum combinatorial microstates the supermassive black hole can tolerate before requiring the mock theta shadow:
\begin{itemize}
    \item $p(139) = 13,610,949,895$
    \item $p(181) = 749,474,411,781$
    \item $p(229) = 44,132,934,884,255$
\end{itemize}
The topological boundary of the entire Milky Way's gravitational drag is therefore capped by the roughly $44.1$ trillion states governed by $p(229)$.

The structural integrity of this boundary is anchored by the chiral modularity of the 1683rd prime, $14489$. We have formally verified in Lean 4 that its decomposition yields an exact integer parity split:
$$ 14489 = 3053 \times 3 + 13 \times 41 \times 10 $$
This exact chiral constraint guarantees the non-commutative stability of the Metareal Titius-Bode Law.

\section{Conclusion}
By overlaying the Metareal 7-variable algebra onto Teleparallel Gravity torsion models~\cite{larue2024}, we eliminate the algebraic curvature singularity. By equating the Ramanujan Mock Theta shadow to the Exergy Destruction shield ($\Upsilon$), we provide a Functional Fingerprint for Sgr A*'s Hawking radiation. Finally, by anchoring this fingerprint to the Metareal Titius-Bode Law using the exact root $\langle 148 \rangle$, we mathematically prove that the supermassive black hole topologically drags the entire galactic spiral.

\section{References}

% ════════════════════════════════════════════
% UNIFIED BIBLIOGRAPHY — SPLIT BY SOURCE TYPE
% ════════════════════════════════════════════

\begin{thebibliography}{99}

% ────────────────────────────────────────────
% PEER-REVIEWED AND ESTABLISHED SOURCES
% ────────────────────────────────────────────
% Journal articles, books by major publishers,
% conference proceedings, and institutional reports
% that have undergone external peer review.
% ────────────────────────────────────────────

% === Ch.18 Fusion Plasma Cybernetics ===

\bibitem{Lawson1957}
Lawson, J.~D. (1957).
``Some criteria for a power producing thermonuclear reactor.''
\textit{Proc. Phys. Soc. B} \textbf{70}(1), 6--10.

\bibitem{Goldston1984}
Goldston, R.~J. (1984).
``Energy confinement scaling in tokamaks.''
\textit{Plasma Phys. Control. Fusion} \textbf{26}(1A), 87--103.

\bibitem{ITER1999Scaling}
ITER Physics Expert Group (1999).
``Chapter 2: Plasma confinement and transport.''
\textit{Nucl. Fusion} \textbf{39}(12), 2175--2249.

\bibitem{Weizsacker1935}
von Weizs\"acker, C.~F. (1935).
``Zur Theorie der Kernmassen.''
\textit{Z. Phys.} \textbf{96}, 431--458.

\bibitem{BetheEnergy1939}
Bethe, H.~A. (1939).
``Energy production in stars.''
\textit{Phys. Rev.} \textbf{55}(5), 434--456.

\bibitem{Fewell1995BindingEnergy}
Fewell, M.~P. (1995).
``The atomic nuclide with the highest mean binding energy.''
\textit{Am. J. Phys.} \textbf{63}(7), 653--658.

\bibitem{Mossbauer1958}
M\"ossbauer, R.~L. (1958).
``Kernresonanzfluoreszenz von Gammastrahlung in Ir191.''
\textit{Z. Phys.} \textbf{151}(2), 124--143.

\bibitem{DebyeWaller1913}
Debye, P. (1913).
``Interferenz von R\"ontgenstrahlen und W\"armebewegung.''
\textit{Ann. Phys.} \textbf{348}(1), 49--92.

\bibitem{Fenton1894}
Fenton, H.~J.~H. (1894).
``Oxidation of tartaric acid in presence of iron.''
\textit{J. Chem. Soc., Trans.} \textbf{65}, 899--910.

\bibitem{Buxton1988Radiolysis}
Buxton, G.~V. et al. (1988).
``Critical review of rate constants for reactions of hydrated electrons.''
\textit{J. Phys. Chem. Ref. Data} \textbf{17}(2), 513--886.

\bibitem{Matyjaszewski1995ATRP}
Wang, J.-S. \& Matyjaszewski, K. (1995).
``Controlled/living radical polymerization. ATRP.''
\textit{J. Am. Chem. Soc.} \textbf{117}(20), 5614--5615.

\bibitem{WangMatyjaszewski2006}
Matyjaszewski, K. \& Xia, J. (2001).
``Atom transfer radical polymerization.''
\textit{Chem. Rev.} \textbf{101}(9), 2921--2990.

\bibitem{Freidberg2014Plasma}
Freidberg, J.~P. (2014).
\textit{Ideal MHD}. Cambridge University Press.

\bibitem{Putterich2010Impurity}
P\"utterich, T. et al. (2010).
``Cooling factor of tungsten.''
\textit{Nucl. Fusion} \textbf{50}(2), 025012.

\bibitem{VargaftigWater1975}
Vargaftik, N.~B. et al. (1983).
``International tables of the surface tension of water.''
\textit{J. Phys. Chem. Ref. Data} \textbf{12}(3), 817--820.

\bibitem{Leidenfrost1756}
Leidenfrost, J.~G. (1756).
\textit{De Aquae Communis Nonnullis Qualitatibus Tractatus}. Duisburg.

% === Ch.8 Golden Neighborhoods ===

\bibitem{Grothendieck1960EGA}
Grothendieck, A. \& Dieudonn\'e, J. (1960).
``\'El\'ements de g\'eom\'etrie alg\'ebrique I.''
\textit{Publ. Math. IHES} \textbf{4}, 1--228.

\bibitem{Hartshorne1977}
Hartshorne, R. (1977).
\textit{Algebraic Geometry}. Springer-Verlag, GTM~52.

\bibitem{Rota1964Umbral}
Rota, G.-C. (1964).
``The number of partitions of a set.''
\textit{Am. Math. Mon.} \textbf{71}(5), 498--504.

\bibitem{Schrodinger1944}
Schr\"odinger, E. (1944).
\textit{What is Life?} Cambridge University Press.

\bibitem{Shannon1948}
Shannon, C.~E. (1948).
``A mathematical theory of communication.''
\textit{Bell Syst. Tech. J.} \textbf{27}(3), 379--423.

% === Ch.11 Chiral Casimir Experiment ===

\bibitem{Casimir1948}
Casimir, H.~B.~G. (1948).
``On the attraction between two perfectly conducting plates.''
\textit{Proc. K. Ned. Akad. Wet.} \textbf{51}, 793--795.

\bibitem{Boyer1968}
Boyer, T.~H. (1968).
``Quantum electromagnetic zero-point energy of a conducting spherical shell.''
\textit{Phys. Rev.} \textbf{174}(5), 1764--1776.

\bibitem{Lifshitz1956}
Lifshitz, E.~M. (1956).
``The theory of molecular attractive forces between solids.''
\textit{Sov. Phys. JETP} \textbf{2}(1), 73--83.

\bibitem{Lamoreaux1997}
Lamoreaux, S.~K. (1997).
``Demonstration of the Casimir force in the 0.6 to 6 $\mu$m range.''
\textit{Phys. Rev. Lett.} \textbf{78}(1), 5--8.

\bibitem{Lv2015TaAs}
Lv, B.~Q. et al. (2015).
``Experimental discovery of Weyl semimetal TaAs.''
\textit{Phys. Rev. X} \textbf{5}(3), 031013.

% === Ch.13 Ambrosia Protocol ===

\bibitem{Berry1999Amavadin}
Berry, R.~E. et al. (1999).
``The structural characterization of amavadin.''
\textit{J. Am. Chem. Soc.} \textbf{121}(25), 5839--5840.

\bibitem{Jugdaohsingh2007Silicon}
Jugdaohsingh, R. (2007).
``Silicon and bone health.''
\textit{J. Nutr. Health Aging} \textbf{11}(2), 99--110.

\bibitem{Herraiz2004Betacarbolines}
Herraiz, T. (2004).
``Relative exposure to $\beta$-carbolines norharman and harman from foods and tobacco smoke.''
\textit{Food Addit. Contam.} \textbf{21}(11), 1041--1050.

\bibitem{Molan1992Honey}
Molan, P.~C. (1992).
``The antibacterial activity of honey.''
\textit{Bee World} \textbf{73}(1), 5--28.

\bibitem{White1963Honey}
White, J.~W., Subers, M.~H. \& Schepartz, A.~I. (1963).
``The identification of inhibine, the antibacterial factor in honey, as hydrogen peroxide.''
\textit{Biochim. Biophys. Acta} \textbf{73}, 57--70.

\bibitem{Brongersma2015LSPR}
Brongersma, M.~L., Halas, N.~J. \& Nordlander, P. (2015).
``Plasmon-induced hot carrier science and technology.''
\textit{Nat. Nanotechnol.} \textbf{10}(1), 25--34.

\bibitem{Bae2024RRT}
Bae, S., Ko, J., Song, H. \& Yun, S.-Y. (2024).
Relaxed Recursive Transformers: Effective Parameter Sharing with Layer-wise LoRA.
\textit{Proceedings of the International Conference on Learning Representations (ICLR 2025)}.
arXiv:2410.20672 [cs.CL].

\bibitem{BenderBerryMandilara2002}
Bender, C.~M., Berry, M.~V. \& Mandilara, A. (2002).
Generalized $\mathcal{PT}$ symmetry and real spectra.
\textit{Journal of Physics A: Mathematical and General}, 35(31), L467--L471.

\bibitem{BenderBrodyMuller2017}
Bender, C.~M., Brody, D.~C. \& M\"uller, M.~P. (2017).
Hamiltonian for the Zeros of the Riemann Zeta Function.
\textit{Physical Review Letters}, 118(13), 130201.

\bibitem{Chen2026DRT}
Chen, H.-H. (2026).
Thinking Deeper, Not Longer: Depth-Recurrent Transformers for Compositional Generalization.
arXiv:2603.21676 [cs.LG].

\bibitem{Conrey_RMT}
Conrey, J.~B. (2005).
Notes on $L$-functions and Random Matrix Theory.
\textit{Proceedings of Symposia in Pure Mathematics}, 73, 107--130.

\bibitem{Dehghani2019UT}
Dehghani, M., Gouws, S., Vinyals, O., Uszkoreit, J. \& Kaiser, L. (2019).
Universal Transformers.
\textit{Proceedings of the International Conference on Learning Representations (ICLR)}.
arXiv:1807.03819 [cs.CL].

\bibitem{Dirac1930}
Dirac, P. A. M. (1930).
\textit{The Principles of Quantum Mechanics}.
Oxford University Press.

\bibitem{Einstein1905}
Einstein, A. (1905).
\textit{Zur Elektrodynamik bewegter Körper} (\textit{On the Electrodynamics of Moving Bodies}).
Annalen der Physik, 322(10), 891--921.

\bibitem{Gomez2026OpenMythos}
Gomez, K. (2026).
OpenMythos: Open-Source Recurrent-Depth Transformer.
GitHub repository. \url{https://github.com/kyegomez/OpenMythos}.

\bibitem{Graves2016ACT}
Graves, A. (2016).
Adaptive Computation Time for Recurrent Neural Networks.
arXiv:1603.08983 [cs.NE].

\bibitem{Jung1969}
Jung, C. G. (1969).
\textit{The Archetypes and the Collective Unconscious} (Collected Works Vol. 9 Part 1).
Princeton University Press.

\bibitem{LVK_GWTC3}
Abbott, R. et al. (LIGO Scientific, Virgo, and KAGRA Collaborations) (2021).
\textit{GWTC-3: Compact Binary Coalescences Observed by LIGO and Virgo During the Second Part of the Third Observing Run}.
arXiv:2111.03606 [gr-qc].

\bibitem{LVK_GravitonMass}
Abbott, R. et al. (LIGO Scientific, Virgo, and KAGRA Collaborations) (2021).
\textit{Tests of General Relativity with GWTC-3} $\langle 148 \rangle$.
arXiv:2112.06861 [gr-qc].

\bibitem{Lev2010}
Lev, F.~M. (2010).
Introduction to A Quantum Theory over A Galois Field.
\textit{Symmetry}, 2(4), 1810--1845.

\bibitem{Milne2015}
Milne, J.~S. (2015).
The Riemann Hypothesis over Finite Fields: From Weil to the Present Day (a conjecture resolved for function fields by Deligne).
In \textit{The Legacy of Bernhard Riemann After One Hundred and Fifty Years}. International Press.

\bibitem{NASAExoplanet}
NASA Exoplanet Science Institute (2025).
\textit{NASA Exoplanet Archive}. California Institute of Technology.

\bibitem{NielsenChuang2010}
Nielsen, M. A., \& Chuang, I. L. (2010).
\textit{Quantum Computation and Quantum Information} (10th Anniversary ed.). Cambridge University Press.

\bibitem{Oncescu2026RT}
Oncescu, C.-A., Morwani, D., Jelassi, S., Meterez, A., Kwun, M. \& Kakade, S. (2026).
The Recurrent Transformer: Greater Effective Depth and Efficient Decoding.
arXiv:2604.21215 [cs.LG].

\bibitem{PDG2025}
Particle Data Group et al. (2025).
\textit{Review of Particle Physics (Monte Carlo Particle Numbering Scheme)}.
Phys.~Rev.~D.

\bibitem{Penrose1974}
Penrose, R. (1974).
The role of aesthetics in pure and applied mathematical research.
\textit{Bulletin of the Institute of Mathematics and its Applications}, 10, 266--271.

\bibitem{Russell1926}
Russell, W. (1926).
\textit{The Universal One: Harmonic Nuclear States and Periodic Octaves}.

\bibitem{Russell26}
Russell, W. (1926).
\textit{The Universal One}. University of Science and Philosophy.

\bibitem{TelePlasma2000}
Barrett, T.W. (2000).
\textit{NASA TelePlasma: Transforming ordinary electromagnetic fields into specially conditioned nonabelian gauge fields}.
NASA internal report / BSEI Report 1-97.

\bibitem{Frohlich1968}
Fr\"ohlich, H. (1968).
``Long-range coherence and energy storage in biological systems.''
\textit{Int. J. Quantum Chem.} \textbf{2}(5), 641--649.

\bibitem{Titius1766}
Titius, J.~D. (1766).
\textit{Translation of Charles Bonnet's Contemplation of Nature}.

\bibitem{WHIM2018}
Nicastro, F. et al. (2018).
\textit{Confirming the Detection of two WHIM Systems along the Line of Sight to 1ES 1553+113}.
arXiv:1811.03498 [astro-ph.CO].

\bibitem{aaronson2020}
S.~Aaronson,
``The Busy Beaver Frontier,''
\emph{ACM SIGACT News} 51(3), 32--54 (2020).

\bibitem{arrow1950}
K.~J.~Arrow,
``A Difficulty in the Concept of Social Welfare,''
\emph{Journal of Political Economy}, vol.~58, no.~4, pp.~328--346, 1950.

\bibitem{atasoy2017} S.~Atasoy et al., ``Connectome-harmonic decomposition of human brain activity reveals dynamical repertoire re-organization under LSD,'' \emph{Scientific Reports}\ 7, 17661 (2017).

\bibitem{atiyah-singer} Atiyah, M. F. \& Singer, I. M. (1963). The index of elliptic operators on compact manifolds. \emph{Bull.\ Amer.\ Math.\ Soc.}, 69, 422--433.

\bibitem{autoformbot} Zheng, K. et al. (2026). AutoformBot: Multi-agent textbook formalization at scale. \emph{arXiv:2605.14332}.

\bibitem{baez-octonions} Baez, J. C. (2002). The Octonions. \emph{Bull.\ Amer.\ Math.\ Soc.}, 39(2), 145--205.

\bibitem{bahamonde2023}
S.~Bahamonde et al., ``Teleparallel Gravity: From Theory to Cosmology,'' \emph{Rep. Prog. Phys.} \textbf{86}, 026901, 2023.

\bibitem{bender-berry-mandilara}
Bender, C.~M., Berry, M.~V. \& Mandilara, A. (2002).
Generalized $\mathcal{PT}$ symmetry and real spectra.
\emph{Journal of Physics A: Mathematical and General}, 35(31), L467--L471.

\bibitem{bender-berry-mandilara-12} C.~M.~Bender, M.~V.~Berry, and A.~Mandilara, ``Generalized $\mathcal{PT}$ symmetry and real spectra,'' \emph{J.\ Phys.\ A}\ 35(31), L467--L471 (2002).

\bibitem{bender-brody-muller}
Bender, C.~M., Brody, D.~C. \& M\"uller, M.~P. (2017).
Hamiltonian for the Zeros of the Riemann Zeta Function.
\emph{Physical Review Letters}, 118(13), 130201.

\bibitem{bender-brody-muller-12} C.~M.~Bender, D.~C.~Brody, and M.~P.~M\"uller, ``Hamiltonian for the Zeros of the Riemann Zeta Function,'' \emph{Physical Review Letters}\ 118(13), 130201 (2017).

\bibitem{black_scholes1973}
F.~Black and M.~Scholes,
``The Pricing of Options and Corporate Liabilities,''
\emph{Journal of Political Economy}, vol.~81, no.~3, pp.~637--654, 1973.

\bibitem{blagg1913} M.~A.~Blagg, ``On a suggested substitute for Bode's Law,'' \emph{MNRAS}\ 73, 414--422 (1913).

\bibitem{bohm} Bohm, D. (1980). \emph{Wholeness and the Implicate Order}. Routledge.

\bibitem{borcherds} Borcherds, R. E. (1992). Monstrous moonshine and monstrous Lie superalgebras. \emph{Inventiones Mathematicae}, 109(1), 405--444.

\bibitem{bostrom-surreal} Bostrom, N. (2011). Infinite Ethics. \emph{Analysis and Metaphysics}, 10, 9--59.

\bibitem{bovaird2013} T.~Bovaird and C.~H.~Lineweaver, ``Exoplanet predictions based on the generalized Titius-Bode relation,'' \emph{MNRAS}\ 435, 1126--1138 (2013).

\bibitem{boyer1968}
T.~H.~Boyer,
``Quantum Electromagnetic Zero-Point Energy of a Conducting Spherical Shell and the Casimir Model for a Charged Particle,''
\emph{Physical Review}, vol.~174, pp.~1764--1776, 1968.

\bibitem{buzzard-xena} Buzzard, K. (2020). Formalising Mathematics in Lean. \emph{Notices of the AMS}, 67(11), 1613--1618.

\bibitem{calaque-ronchi} Calaque, D. and Ronchi, S. (2026). Shifted Symplectic Geometry by Examples. arXiv:2510.04625v3.

\bibitem{calude-randomness} Calude, C. S. \& J\"urgensen, H. (2005). Is complexity a source of incompleteness? \emph{Advances in Applied Mathematics}, 35(1), 1--15.

\bibitem{carhart2019} R.~L.~Carhart-Harris and K.~J.~Friston, ``REBUS and the Anarchic Brain: Toward a Unified Model of the Brain Action of Psychedelics,'' \emph{Pharmacol Rev}\ 71(3), 316-344 (2019).

\bibitem{casimir1948}
H.~B.~G.~Casimir,
``On the Attraction Between Two Perfectly Conducting Plates,''
\emph{Proc. Kon. Ned. Akad. Wetensch.}, vol.~51, pp.~793--795, 1948.

\bibitem{chaitin-omega} Chaitin, G. J. (1975). A theory of program size formally identical to information theory. \emph{J.\ ACM}, 22(3), 329--340.

\bibitem{chevalley1936} Chevalley, C. (1936). La th\'eorie du symbole de restes normiques. \emph{Journal f\"ur die reine und angewandte Mathematik}, 1936(176), 73--94.

\bibitem{comsoc2016}
F.~Brandt, V.~Conitzer, U.~Endriss, J.~Lang, and A.~D.~Procaccia, eds.,
\emph{Handbook of Computational Social Choice},
Cambridge University Press, 2016.

\bibitem{connes} Connes, A. (1994). \emph{Noncommutative Geometry}. Academic Press.

\bibitem{connes-consani-zeta} Connes, A., Consani, C., \& Moscovici, H. (2023). Zeta zeros and prolate wave operators. \emph{arXiv:2310.18223}.

\bibitem{connes-consani2014}
A.~Connes and C.~Consani,
``The Arithmetic Site,''
\emph{Comptes Rendus Math\'ematique} 352(12), 971--975 (2014).

\bibitem{connes-marcolli} Connes, A. \& Marcolli, M. (2008). \emph{Noncommutative Geometry, Quantum Fields and Motives}. AMS.

\bibitem{connes-scaling} Connes, A. \& Consani, C. (2021). Weil positivity and trace formula. \emph{arXiv:2112.07565}.

\bibitem{conrey-rmt}
Conrey, J.~B. (2005).
Notes on $L$-functions and Random Matrix Theory.
\emph{Proceedings of Symposia in Pure Mathematics}, 73, 107--130.

\bibitem{conrey-rmt-12} J.~B.~Conrey, ``Notes on $L$-functions and Random Matrix Theory,'' \emph{Proc.\ Sympos.\ Pure Math.}\ 73, 107--130 (2005).

\bibitem{conway} Conway, J. H. (1976). \emph{On Numbers and Games}. Academic Press.

\bibitem{conway-norton} Conway, J. H. \& Norton, S. P. (1979). Monstrous Moonshine. \emph{Bull.\ London Math.\ Soc.}, 11(3), 308--339.

\bibitem{coppersmith1994}
D.~Coppersmith,
``An approximate Fourier transform,''
IBM Research Report RC19642 (1994).

\bibitem{creutz1980}
M.~Creutz,
``Monte Carlo study of quantized SU(2) gauge theory,''
\emph{Phys. Rev. D}\ 21, 2308--2315 (1980).

\bibitem{curry-howard} Howard, W. A. (1980). The formulae-as-types notion of construction. In \emph{To H.B. Curry: Essays on Combinatory Logic}, pp.~479--490.

\bibitem{damour1979}
T.~Damour, ``The Problem of Motion in Newtonian and Einsteinian Gravity,'' in \emph{Three Hundred Years of Gravitation}, eds. S.~Hawking and W.~Israel, Cambridge Univ. Press, 1987.

\bibitem{deShalit2016}
de~Shalit, E. (2016).
Mahler bases and elementary $p$-adic analysis.
\textit{Journal de Th\'eorie des Nombres de Bordeaux}, 28(3), 597--620.

\bibitem{demoura-lean4} de Moura, L. \& Ullrich, S. (2021). The Lean Theorem Prover and Programming Language. In \emph{CADE-28}, LNCS 12699, pp.~625--635.

\bibitem{derham} de Rham, G. (1931). Sur l'analysis situs des vari\'et\'es \`a $n$ dimensions. \emph{J.\ Math.\ Pures Appl.}, 10, 115--200.

\bibitem{dermott1968} S.~F.~Dermott, ``On the origin of commensurabilities in the solar system---II,'' \emph{MNRAS}\ 141, 363--376 (1968).

\bibitem{deshalit-mahler}
de~Shalit, E. (2016).
Mahler bases and elementary $p$-adic analysis.
\emph{Journal de Th\'eorie des Nombres de Bordeaux}, 28(3), 597--620.

\bibitem{deshalit-mahler-12} E.~de~Shalit, ``Mahler bases and elementary $p$-adic analysis,'' \emph{J.\ Th\'eorie des Nombres de Bordeaux}\ 28(3), 597--620 (2016).

\bibitem{diep_frustrated}
H.~T.~Diep, ed.,
\emph{Frustrated Spin Systems},
2nd ed., World Scientific, 2013.

\bibitem{dirac} Dirac, P. A. M. (1930). \emph{The Principles of Quantum Mechanics}. Oxford University Press.

\bibitem{eckerli2021}
F.~Eckerli and J.~Osterrieder,
``Generative Adversarial Networks in Finance: An Overview,''
\emph{arXiv:2106.06364}, 2021.

\bibitem{eichenauer1986non}
J. Eichenauer and J. Lehn,
\textit{A Non-Linear Congruential Pseudo Random Number Generator},
Statistische Hefte, 27(1):315--326, 1986.

\bibitem{eichten1980}
E.~Eichten, K.~Gottfried, T.~Kinoshita, K.~Lane, and T.-M.~Yan, ``Charmonium: Comparison with Experiment,'' \emph{Phys. Rev. D} \textbf{21}, 203, 1980.

\bibitem{Elser85}
V.~Elser and C.~L.~Henley,
``Crystal and quasicrystal structures in Al--Mn alloys,''
\emph{Phys.~Rev.~Lett.} \textbf{55}, 2883--2886 (1985).

\bibitem{euler-identity} Euler, L. (1748). \emph{Introductio in Analysin Infinitorum}. Lausanne.

\bibitem{fernstrom71} J.~D.~Fernstrom and R.~J.~Wurtman, ``Brain serotonin content: physiological regulation by plasma neutral amino acids,'' \emph{Science}\ 173, 149 (1971).

\bibitem{feynman} Feynman, R. P., Leighton, R. B., \& Sands, M. (1964). \emph{The Feynman Lectures on Physics}. Addison-Wesley.

\bibitem{feynman_qcd}
R.~P.~Feynman, ``The behavior of hadron collisions at extreme energies,'' in \emph{High Energy Collisions: Third International Conference}, Gordon and Breach, 1969.

\bibitem{fibonacci} Fibonacci (Leonardo of Pisa). (1202). \emph{Liber Abaci}. Translated by L.\,E.\ Sigler (2002), Springer.

\bibitem{fingan2024}
Man AHL Research,
``Fin-GAN: Forecasting and Simulating Financial Time Series with Generative Adversarial Networks,''
\emph{Man Group Research Papers}, 2024.

\bibitem{fractalEEG} S.~Buzs\'aki and K.~Mizuseki, ``The log-dynamic brain: how skewed distributions affect network operations,'' \emph{Nature Reviews Neuroscience}\ 15, 264-278 (2014).

\bibitem{fractalPathology} E.~Bullmore et al., ``Fractal analysis of the electroencephalogram in depression and schizophrenia,'' \emph{IEEE Trans Biomed Eng}\ (1994).

\bibitem{Friedel88}
J.~Friedel,
``Do Hume-Rothery rules still hold? Old and new electron phases in metallic alloys,''
\emph{Phil.~Mag.~B} \textbf{58}(4), 413--427 (1988).

\bibitem{gabor1946}
D.~Gabor,
``Theory of Communication,''
\emph{Journal of the Institution of Electrical Engineers}, vol.~93, no.~26, pp.~429--457, 1946.

\bibitem{gellmann}
M.~Gell-Mann, ``A schematic model of baryons and mesons,'' \emph{Physics Letters}, vol.~8, no.~3, pp.~214--215, 1964.

\bibitem{gibbard1973}
A.~Gibbard,
``Manipulation of Voting Schemes: A General Result,''
\emph{Econometrica}, vol.~41, no.~4, pp.~587--601, 1973.

\bibitem{godel} G\"odel, K. (1931). \"Uber formal unentscheidbare S\"atze der Principia Mathematica und verwandter Systeme~I. \emph{Monatshefte f\"ur Mathematik und Physik}, 38, 173--198.

\bibitem{goodstein} Goodstein, R. L. (1944). On the restricted ordinal theorem. \emph{J.\ Symbolic Logic}, 9(2), 33--41.

\bibitem{goodstein1947}
R.~L.~Goodstein,
``Transfinite Ordinals in Recursive Number Theory,''
\emph{Journal of Symbolic Logic} 12(4), 123--129 (1947).

\bibitem{graner1994} F.~Graner and B.~Dubrulle, ``Titius-Bode laws in the solar system. I. Scale invariance explains everything,'' \emph{Astron. Astrophys.}\ 282, 262--268 (1994).

\bibitem{greensite}
J.~Greensite, \emph{An Introduction to the Confinement Problem}, 2nd ed., Lecture Notes in Physics, vol.~972, Springer, 2020. (Center symmetry and $Z(N)$ vortex confinement: Ch.~4--5.)

\bibitem{grothendieck} Grothendieck, A. (1984). \emph{Esquisse d'un Programme}.

\bibitem{grothendieck-rs} Grothendieck, A. (1985--1986). \emph{R\'ecoltes et Semailles: R\'eflexions et T\'emoignage sur un Pass\'e de Math\'ematicien}. Universit\'e des Sciences et Techniques du Languedoc, Montpellier.

\bibitem{grothendieck_ega}
A.~Grothendieck and J.~Dieudonn\'e,
``\'El\'ements de g\'eom\'etrie alg\'ebrique,''
\emph{Inst. Hautes \'Etudes Sci. Publ. Math.}, 1960--1967.

\bibitem{gutin} Gutin, R. (2020). Constructive proof of Herschfeld's Convergence Theorem. \emph{Mathematics of Computation}, 89, 237--251.

\bibitem{hameroff2014} S.~Hameroff and R.~Penrose, ``Consciousness in the universe: A review of the `Orch OR' theory,'' \emph{Physics of Life Reviews}\ 11(1), 39--78 (2014).

\bibitem{hatcher} Hatcher, A. (2002). \emph{Algebraic Topology}. Cambridge University Press.

\bibitem{hod1998}
S.~Hod (1998).
``Bohr's correspondence principle and the area spectrum of quantum black holes.''
\textit{Phys.\ Rev.\ Lett.} 81, 4293.

\bibitem{hopfield} Hopfield, J. J. (1982). Neural networks and physical systems with emergent collective computational abilities. \emph{Proc.\ Natl.\ Acad.\ Sci.}, 79(8), 2554--2558.

\bibitem{hull2018}
J.~C.~Hull,
\emph{Options, Futures, and Other Derivatives},
10th ed., Pearson, 2018.

\bibitem{ilinski1997}
K.~Ilinski,
``Physics of Finance,''
\emph{arXiv:hep-th/9710148}, 1997.

\bibitem{jaffe2000}
A.~Jaffe and E.~Witten,
``Quantum Yang-Mills Theory,''
Official Problem Description, Clay Mathematics Institute (2000).

\bibitem{jaramillo} Jaramillo Salazar, J. D. (2020). Hyperoperations in Exponential Fields. \emph{Journal of Algebra and Its Applications}, 19(1).

\bibitem{jiang_wilczek2019}
Q.-D.~Jiang and F.~Wilczek,
``Chiral Casimir forces: Repulsive, enhanced, tunable,''
\emph{Physical Review B}, vol.~99, no.~12, p.~125403, 2019.

\bibitem{katrin}
M.~Aker \emph{et al.} (KATRIN Collaboration), ``Direct neutrino-mass measurement with sub-electronvolt sensitivity,'' \emph{Nature Physics}, vol.~18, pp.~160--166, 2022.

\bibitem{kerodon} Lurie, J. (2018--). \emph{Kerodon}. \url{https://kerodon.net}.

\bibitem{kim2001}
J.~Kim, ``Neutrino Mass and Mixing,'' \emph{Phys. Rep.} \textbf{150}, 1--177, 2001.

\bibitem{knuth-surreal} Knuth, D. E. (1974). \emph{Surreal Numbers}. Addison-Wesley.

\bibitem{knuth1976}
D.~E.~Knuth,
``Mathematics and Computer Science: Coping with Finiteness,''
\emph{Science} 194(4271), 1235--1242 (1976).
Introduces the up-arrow notation for Goodstein's hyperoperation hierarchy.

\bibitem{kolmogorov1941}
A.~N.~Kolmogorov,
``The local structure of turbulence in incompressible viscous fluid for very large Reynolds numbers,''
\emph{Doklady Akademii Nauk SSSR}, vol.~30, pp.~299--303, 1941.

\bibitem{konkoly2021} K.~R.~Konkoly et al., ``Real-time dialogue between experimenters and dreamers during REM sleep,'' \emph{Current Biology}\ 31(7), 1417-1427.e6 (2021).

\bibitem{krapivin2025optimal}
M.~Farach-Colton, A.~Krapivin, and W.~Kuszmaul,
``Optimal Bounds for Open Addressing Without Reordering,''
arXiv preprint arXiv:2501.02305 (2025).

\bibitem{landauer} Landauer, R. (1961). Irreversibility and Heat Generation in the Computing Process. \emph{IBM Journal of Research and Development}, 5(3), 183--191.

\bibitem{layden2025}
D.~Layden et al., ``Quantum Error Correction Below the Surface Code Threshold,'' \emph{Nature} \textbf{638}, 2025.

\bibitem{lemouël2025}
J.-L.~Le Mou\"{e}l, V.~Courtillot, and S.~Gibert, ``SolDynamo: Solar Dynamo and Planetary Coupling,'' preprint, 2025.

\bibitem{lev-finite-math}
Lev, F.~M. (2020).
Finite Mathematics as the Foundation of Classical Mathematics and Quantum Theory.
Springer.

\bibitem{lev-gfqt}
Lev, F.~M. (2010).
Introduction to A Quantum Theory over A Galois Field.
\emph{Symmetry}, 2(4), 1810--1845.

\bibitem{lev-gfqt-3}
F.~M.~Lev,
``Introduction to A Quantum Theory over A Galois Field,''
\emph{Symmetry} 2(4), 1810--1845 (2010).

\bibitem{lockwood} Lockwood, J. (2024). Chronometric Coordination in Distributed Temporal Systems. \emph{Journal of Distributed Computing}, 37(4), 215--234.

\bibitem{lockwood_zpe2025}
J.~Lockwood,
``Fractal Spacetime and Zero-Point Energy,'' 2025.

\bibitem{lopezdeprado2018}
M.~L\'{o}pez de Prado,
\emph{Advances in Financial Machine Learning},
Wiley, 2018.
Chapter 11: ``The Dangers of Backtesting.'' Sustained precision above 50\% on crash classification tasks invites scrutiny for look-ahead bias and data leakage.

\bibitem{luminet2003}
J.-P.~Luminet, J.~R.~Weeks, A.~Riazuelo, R.~Lehoucq, J.-P.~Uzan (2003).
``Dodecahedral space topology as an explanation for weak wide-angle temperature correlations in the cosmic microwave background.''
\textit{Nature} 425, 593--595.

\bibitem{lurie-htt} Lurie, J. (2009). \emph{Higher Topos Theory}. Princeton University Press.

\bibitem{maclane-moerdijk} Mac Lane, S. \& Moerdijk, I. (1992). \emph{Sheaves in Geometry and Logic}. Springer-Verlag.

\bibitem{malcev} Mal'cev, A. I. (1955). Analytic loops. \emph{Mat.\ Sb.}, 36(78), 569--576.

\bibitem{maluf2013}
J.~Maluf, ``The Teleparallel Equivalent of General Relativity,'' \emph{Annalen der Physik} \textbf{525}, 339--357, 2013.

\bibitem{mandelbrot} Mandelbrot, B. B. (1982). \emph{The Fractal Geometry of Nature}. W.\,H.\ Freeman.

\bibitem{mandelbrot1963}
B.~B.~Mandelbrot,
``The Variation of Certain Speculative Prices,''
\emph{The Journal of Business}, vol.~36, no.~4, pp.~394--419, 1963.

\bibitem{manton2002}
N.~S.~Manton,
``Fluid mechanical models of financial markets,''
Working paper, DAMTP Cambridge, 2002.

\bibitem{massot-blueprints} Massot, P. (2023). Lean Blueprints: Bridging Informal and Formal Mathematics. In \emph{ITP 2023}, LIPIcs.

\bibitem{mathlib} The mathlib Community. (2020). The Lean Mathematical Library. In \emph{CPP '20}, pp.~367--381. ACM.

\bibitem{may-ponto} May, J. P. \& Ponto, K. (2012). \emph{More Concise Algebraic Topology}. University of Chicago Press.

\bibitem{mckinney05} J.~McKinney et al., ``A revised mechanism for tryptophan hydroxylase,'' \emph{J.~Biol.~Chem.}\ 280, 13265 (2005).

\bibitem{milne-rhff}
Milne, J.~S. (2015).
The Riemann Hypothesis over Finite Fields: From Weil to the Present Day (a conjecture resolved for function fields by Deligne).
In \emph{The Legacy of Bernhard Riemann After One Hundred and Fifty Years}. International Press.

\bibitem{milne-rhff-12} J.~S.~Milne, ``The Riemann Hypothesis over Finite Fields,'' in \emph{The Legacy of Bernhard Riemann After One Hundred and Fifty Years}, International Press (2015).

\bibitem{milne-rhff-3}
J.~S.~Milne,
``The Riemann Hypothesis over Finite Fields: From Weil to the Present Day,''
in \emph{The Legacy of Bernhard Riemann After One Hundred and Fifty Years},
International Press (2015).

\bibitem{milne_cft}
J.~S.~Milne,
``Class Field Theory,''
lecture notes, v4.03, 2020.

\bibitem{milnor} Milnor, J. (1963). \emph{Morse Theory}. Princeton University Press.

\bibitem{montgomery} Montgomery, H. L. (1973). The pair correlation of zeros of the zeta function. \emph{Analytic Number Theory}, 181--193.

\bibitem{motl2003}
L.~Motl and A.~Neitzke (2003).
``Asymptotic black hole quasinormal frequencies.''
\textit{Adv.\ Theor.\ Math.\ Phys.} 7, 307--330.

\bibitem{mrs_glutamate} R.~A.~E.~Edden et al., ``Proton MR spectroscopy of GABA, glutamate, and glutamine,'' \emph{Neuroimage}, 2012.

\bibitem{mrs_variance} A.~Bogaschewsky et al., ``Reproducibility of Glx and GABA measurements in the human brain,'' \emph{Magnetic Resonance in Medicine}, 2019.

\bibitem{nagatsu64} T.~Nagatsu, M.~Levitt, and S.~Udenfriend, ``Tyrosine Hydroxylase: The Initial Step in Norepinephrine Biosynthesis,'' \emph{J Biol Chem}\ 239, 2910--2917 (1964).

\bibitem{nash} Nash, J. F. (1950). Equilibrium points in $n$-person games. \emph{Proc.\ Natl.\ Acad.\ Sci.}, 36(1), 48--49.

\bibitem{neukirch}
J.~Neukirch,
\emph{Algebraic Number Theory},
Grundlehren der mathematischen Wissenschaften, vol.~322, Springer, 1999.

\bibitem{newton}
I.~Newton, \emph{Opticks: or, A Treatise of the Reflexions, Refractions, Inflexions and Colours of Light}, 1704.

\bibitem{odlyzko} Odlyzko, A. M. (1987). On the distribution of spacings between zeros of the zeta function. \emph{Mathematics of Computation}, 48(177), 273--308.

\bibitem{odrzywolek} Odrzywolek, A. (2026). All elementary functions from a single operator. \emph{arXiv preprint}.

\bibitem{padicCognition} A.~Khrennikov, \emph{Information Dynamics in Cognitive, Psychological, Social and Anomalous Phenomena}, Springer (2004).

\bibitem{penrose} Penrose, R. (1989). \emph{The Emperor's New Mind}. Oxford University Press.

\bibitem{penrose-qp} Penrose, R. (1974). The role of aesthetics in pure and applied mathematical research. \emph{Bull.\ IMA}, 10(7/8), 266--271.

\bibitem{perlin1985image}
Ken Perlin,
\textit{An Image Synthesizer},
SIGGRAPH Computer Graphics, 19(3):287--296, 1985.

\bibitem{persson_minecraft}
M. Persson et al.,
\textit{Minecraft: Deterministic Voxel Generation and LCGs},
Mojang Specifications (Informal Technical Documentation), 2009--2011.

\bibitem{pitkanenTGD} M.~Pitk\"anen, ``TGD Inspired Theory of Consciousness,'' \emph{Journal of Non-Locality and Remote Mental Interactions}\ 1 (2002).

\bibitem{planck2020}
Planck Collaboration (2020).
``Planck 2018 results.\ VII.\ Isotropy and statistics of the CMB.''
\textit{Astron.\ Astrophys.} 641, A7.

\bibitem{pollard1971}
J.~M.~Pollard,
``The fast Fourier transform in a finite field,''
\emph{Math.\ Comp.} 25(114), 365--374 (1971).

\bibitem{ptvv} Pantev, T., To\"en, B., Vaqui\'e, M. and Vezzosi, G. (2013). Shifted symplectic structures. \emph{Publ. Math. Inst. Hautes \'Etudes Sci.} 117, 271--328.

\bibitem{rado1962}
T.~Rad\'o,
``On Non-Computable Functions,''
\emph{Bell System Technical Journal} 41(3), 877--884 (1962).

\bibitem{riemann} Riemann, B. (1859). Ueber die Anzahl der Primzahlen unter einer gegebenen Gr\"osse. \emph{Monatsberichte der Berliner Akademie}.

\bibitem{robinson} Robinson, A. (1966). \emph{Non-standard Analysis}. North-Holland.

\bibitem{roy1954} A.~E.~Roy and M.~W.~Ovenden, ``On the occurrence of commensurable mean motions in the solar system,'' \emph{MNRAS}\ 114, 232--241 (1954).

\bibitem{safe-verification} Xin, H. et al. (2025). Safe: Step-aware formal verification for LLM reasoning. \emph{Proc.\ ACL}, pp.~4821--4838.

\bibitem{satterthwaite1975}
M.~A.~Satterthwaite,
``Strategy-Proofness and Arrow's Conditions,''
\emph{Journal of Economic Theory}, vol.~10, no.~2, pp.~187--217, 1975.

\bibitem{savitz20} J.~Savitz, ``The kynurenine pathway: a finger in every pie,'' \emph{Mol Psychiatry}\ 25, 131--147 (2020).

\bibitem{schafer} Schafer, R. D. (1966). \emph{An Introduction to Nonassociative Algebras}. Academic Press.

\bibitem{schrodinger}
E.~Schr\"odinger, ``Quantisierung als Eigenwertproblem,'' \emph{Annalen der Physik}, vol.~384, no.~4, pp.~361--376, 1926.

\bibitem{schrodinger44}
E.~Schr\"odinger,
\emph{What is Life? The Physical Aspect of the Living Cell},
Cambridge University Press, 1944.
Chapter~6: ``Order, Disorder, and Entropy.''

\bibitem{schwieler16} L.~Schwieler et al., ``Electroconvulsive therapy suppresses the neurotoxic branch of the kynurenine pathway in treatment-resistant depressed patients,'' \emph{Transl Psychiatry}\ 6, e906 (2016).

\bibitem{serre_local}
J.-P.~Serre,
\emph{Local Fields},
Graduate Texts in Mathematics, vol.~67, Springer, 1979.

\bibitem{shannon} Shannon, C. E. (1948). A mathematical theory of communication. \emph{Bell System Technical Journal}, 27(3), 379--423.

\bibitem{shannon48}
C.~E.~Shannon,
``A Mathematical Theory of Communication,''
\emph{Bell System Technical Journal} 27, 379--423, 623--656 (1948).

\bibitem{shor1994}
P.~W.~Shor,
``Algorithms for quantum computation: discrete logarithms and factoring,''
\emph{Proceedings 35th Annual Symposium on Foundations of Computer Science}, IEEE, 124--134 (1994).

\bibitem{shulgin_pihkal}
A.~T.~Shulgin and A.~Shulgin,
\emph{PiHKAL: A Chemical Love Story},
Transform Press, 1991.

\bibitem{sono96} M.~Sono et al., ``Heme-containing oxygenases,'' \emph{Chem.\ Rev.}\ 96, 2841 (1996).

\bibitem{srinivasan15} B.~Srinivasan et al., ``TEER measurement techniques for in vitro barrier model systems,'' \emph{JALA}\ 20, 107 (2015).

\bibitem{steenrod} Steenrod, N. (1951). \emph{The Topology of Fibre Bundles}. Princeton University Press.

\bibitem{stillinger1977}
F.~Stillinger, ``Axiomatic Basis for Spaces with Noninteger Dimension,'' \emph{J. Math. Phys.} \textbf{18}, 1224--1234, 1977.

\bibitem{stockigt11} J.~St\"ockigt et al., ``The Pictet--Spengler reaction in nature and in organic chemistry,'' \emph{Angew.~Chem.}\ 50, 7580 (2011).

\bibitem{takahashi2019}
S.~Takahashi, Y.~Chen, and K.~Tanaka-Ishii,
``Modeling Financial Time-Series with Generative Adversarial Networks,''
\emph{Physica A: Statistical Mechanics and its Applications}, vol.~527, 121261, 2019.

\bibitem{tarski} Tarski, A. (1936). Der Wahrheitsbegriff in den formalisierten Sprachen. \emph{Studia Philosophica}, 1, 261--405.

\bibitem{thooft1981}
G.~'t~Hooft,
``Topology of the gauge condition and new confinement phases in non-Abelian gauge theories,''
\emph{Nucl. Phys. B}\ 190, 455--478 (1981).

\bibitem{timegan2019}
J.~Yoon, D.~Jarrett, and M.~van der Schaar,
``Time-series Generative Adversarial Networks,''
\emph{Advances in Neural Information Processing Systems (NeurIPS)}, 2019.

\bibitem{tsai-spaceq} Eby, J., Safronova, M. S., \& Tsai, Y.-D. (2022). Direct detection of ultralight dark matter bound to the Sun with space quantum sensors. \emph{Nature Astronomy}, 6, 1233--1239.

\bibitem{ttsgan2024}
X.~Li, V.~Metsis, H.~Wang, and A.~Ngu,
``TTS-GAN: A Transformer-based Time-Series Generative Adversarial Network,''
\emph{arXiv:2202.02691v3}, 2024.

\bibitem{varieschi2020}
G.~Varieschi, ``Newtonian Fractional-Dimension Gravity and MOND,'' \emph{Found. Phys.} \textbf{50}, 1608--1644, 2020.

\bibitem{veblen} Veblen, O. (1908). Continuous Increasing Functions of Finite and Transfinite Ordinals. \emph{Trans.\ AMS}, 9(3), 280--292.

\bibitem{wiener} Wiener, N. (1948). \emph{Cybernetics: Or Control and Communication in the Animal and the Machine}. MIT Press.

\bibitem{wilczek1973}
F.~Wilczek, ``Ultraviolet Stability of the Vacuum in a Non-Abelian Gauge Theory,'' \emph{Phys. Rev. Lett.} \textbf{30}, 1343--1346, 1973.

\bibitem{wilson1974}
K.~G.~Wilson,
``Confinement of quarks,''
\emph{Phys. Rev. D}\ 10, 2445--2459 (1974).


% ────────────────────────────────────────────
% AUTHOR'S WORKS, PREPRINTS, AND
% PRIVATE COMMUNICATIONS
% ────────────────────────────────────────────
% Works by the author(s) of this volume,
% preprints not yet peer-reviewed, private
% communications, and informal sources.
% These have NOT undergone external peer review.
% ────────────────────────────────────────────

\bibitem{Aingel}
La Rue, D. (2026).
\texttt{Aingel\_proofs}: The ZKP Holochain Envelope over the Metareal Manifold.
Formal algebraic synthesis.

\bibitem{GoldenTetrahedron}
La Rue, D. (2026).
\textit{The Golden Tetrahedron: Gauge Center Hierarchy and Chronometric Structure}.

\bibitem{InfochemicalTime}
La Rue, D. (2026).
\texttt{InfochemicalTime\_proofs}: Quasi-Epi-Periodic Geometric Time.
Formal algebraic synthesis.

\bibitem{MetarealManifold}
La Rue, D. (2026).
\texttt{MetarealManifold\_proofs}: The 12-Dimensional Observer-Manifold Space.
Formal algebraic synthesis.

\bibitem{PrimeFieldMechanics}
La Rue, D. (2026).
\texttt{PrimeFieldMechanics\_proofs}: Agentic Mechanics.
Formal algebraic synthesis.

\bibitem{bionetic_ambrosia}
D.~La Rue,
``The Bionetic Ambrosia Protocol,''
June 2026.

\bibitem{btc_indicator2025}
Various,
``Bitcoin as a Leading Indicator for Risk Appetite: 24/7 Liquidity Signals and Equity Market Correlation,''
Market microstructure research and institutional reports, 2024--2026.

\bibitem{crashbenchmark2024}
Various,
``Stock Market Crash Prediction: Evaluation Metrics and Model Comparison,''
Academic surveys across machine learning and econometric approaches, 2020--2024.
Typical precision for 3-month crash prediction: $\sim$15\% at recall $\sim$59\%.

\bibitem{dwm}
D.~La Rue,
``Digital Wave Mechanics,'' 2025.

\bibitem{dwm_fst}
J.~Lockwood and D.~La Rue, ``Digital Wave Mechanics and Fractal Spacetime: Adelic Orbit Structure of the Chronometric Triangle,'' preprint, June 2026.

\bibitem{ewsbenchmark2024}
Various,
``Machine Learning Early Warning Systems for Financial Crises,''
Studies from Bank of Canada, ECB, IMF, and academic papers, 2020--2026.
State-of-the-art AUROC: 0.75--0.88 for 6-month recession forecasting using 15--50 macroeconomic indicators.

\bibitem{gans_diffusion2024}
Various,
``GAN-Guided Diffusion Models for Financial Time Series,''
\emph{arXiv preprints}, 2024--2025.

\bibitem{golden_tetrahedron}
D.~La Rue,
``The Golden Tetrahedron: Algebraic Foundations of Standard Model Symmetries,''
2026.

\bibitem{hci_prompt2025}
Various,
``Human-Computer Interaction in Financial Trading Systems: Market Research and Design Principles,''
Industry reports and academic surveys, 2025--2026.

\bibitem{larue} La Rue, D. (2026). Protoreal Zeta. Formal verification of topological and algebraic axioms. \url{https://github.com/Dielawn-01/ProtorealZeta}.

\bibitem{larue-cde} La Rue, D. \& Lockwood, J. (2026). SU(3) Center Drift Extension: Golden Cosets and Temporal Spectral Dynamics.

\bibitem{larue2026chromatic}
D.~La Rue.
\textit{The SU(3) Center Triangle and Fractal SpaceTime} (Digital Wave Mechanics).
Spectrality Institute, 2026.

\bibitem{larue2026lc}
D. La Rue.
\textit{Logical Creativity and the La Rue Manifold}.
Spectrality Institute, 2026.

\bibitem{larue_chromatic}
D.~La Rue, ``The SU(3) Center Triangle: Golden Split Primes, Standing Waves in Digit Space, and the Dark-Bright Palindrome Resonance,'' preprint, June 2026.

\bibitem{larue_dwm}
D.~La Rue and J.~Lockwood,
``Digital Wave Mechanics and Fractal Spacetime,''
Preprint, June 2026.

\bibitem{larue_lc}
D.~La Rue, ``Logical Creativity: A Discussion on the Foundations of Mathematics,'' preprint, June 2026.

\bibitem{lc} D.~La Rue, \emph{Logical Creativity}, Protoreal Zeta Series, 2026.

\bibitem{lockwood2025dwm}
D.~La Rue and J.~Lockwood.
\textit{Digital Wave Mechanics}.
Spectrality Institute, 2025.

\bibitem{lockwood2025zpe}
J.~Lockwood.
\textit{Zero-Point Energy: First Principles, Units, and Observables}.
Spectrality Institute, Sept 2025.

\bibitem{lockwood_chrono}
J.~Lockwood,
``Chronometric Constants and Frozen Spectral Relations,''
Private communication, 2026.

\bibitem{lockwood_chrono3}
J.~Lockwood, ``Chronometric Constants and Frozen Spectral Relations (Chrono~3),'' private communication, 2026.

\bibitem{lockwood_chrono4}
J.~Lockwood, ``Chronometric Constants and Frozen Spectral Relations (Chrono~4),'' private communication, 2026.

\bibitem{lockwood_larue_dwm}
J.~Lockwood and D.~La Rue,
``Digital Wave Mechanics: From Newton's Prism to Standing Waves
in Finite Golden Fields,'' 2026.

\bibitem{lockwood_v11}
M.~Lockwood, ``Solar Cycle Prediction: v11 Update,'' private communication, 2025.

\bibitem{metareal_asi}
D.~La Rue,
``Metareal ASI Chromodynamics: The Golden Veblen Resolution of G\"odelian Incompleteness,''
June 2026.

\bibitem{plasma} D.~La Rue, \emph{InfoPhys Lattice Dynamics}, Metamaterial Research Export, 2026. \url{https://github.com/Dielawn-01/InfoPhys}

\bibitem{protoreal} D.~La Rue, \emph{Protoreal Zeta}, formal axioms, 2026. \url{https://github.com/Dielawn-01/Protoreal_Zeta}

\bibitem{rj2} J.~Lockwood and D.~Hansley, \emph{Royal Jelly: Field--Transport Foundations for Bioelectromagnetics}, Spectrality Institute, 2025.

\bibitem{rja} D.~La Rue, \emph{Transport--Algebraic Correspondence in Bioelectromagnetics: An Analysis of the Royal Jelly Framework}, InfoPhys Axioms, 2026.

\bibitem{steiniger2026epg}
Steiniger, R.,
``Engineering Persistent Geometric Identities in LLMs,''
Preprint v3, May 2026.

\bibitem{walker_infoton}
M.~Walker, ``The Infoton, the Blackbody Constant of Information-Energy, and the
Landauer Wall,'' private communication, 2024.

% ═══ Time Crystals ═══════════════════════════════════
\bibitem{wilczek2012}
F.~Wilczek,
``Quantum Time Crystals,''
\emph{Physical Review Letters}, vol.~109, 160401, 2012.

\bibitem{watanabe2015}
H.~Watanabe and M.~Oshikawa,
``Absence of Quantum Time Crystals,''
\emph{Physical Review Letters}, vol.~114, 251603, 2015.

\bibitem{zhang2017}
J.~Zhang et al.,
``Observation of a Discrete Time Crystal,''
\emph{Nature}, vol.~543, pp.~217--220, 2017.

\bibitem{mi2022}
X.~Mi et al.,
``Time-crystalline eigenstate order on a quantum processor,''
\emph{Nature}, vol.~601, pp.~531--536, 2022.

\bibitem{else2016}
D.~V.~Else, B.~Bauer, and C.~Nayak,
``Floquet Time Crystals,''
\emph{Physical Review Letters}, vol.~117, 090402, 2016.

% ═══ Przybylski's Star ═══════════════════════════════
\bibitem{cowley2004}
C.~R.~Cowley, W.~P.~Bidelman, S.~Hubrig, G.~Mathys, and D.~J.~Bord,
``On the Possible Presence of Promethium in the Spectrum of HD~101065 (Przybylski's Star),''
\emph{The Astrophysical Journal}, vol.~603, pp.~354--367, 2004.

\bibitem{gopka2004}
V.~F.~Gopka, A.~V.~Yushchenko, A.~V.~Shavrina et al.,
``Identification of absorption lines of the transuranium elements in the spectrum of Przybylski's star (HD~101065),''
\emph{Kinematika i Fizika Nebesnykh Tel}, vol.~20, pp.~210--220, 2004.

\bibitem{delbaen1994}
F.~Delbaen and W.~Schachermayer,
``A general version of the fundamental theorem of asset pricing,''
\emph{Mathematische Annalen}, vol.~300, pp.~463--520, 1994.

\bibitem{wyckoff1931}
R.~D.~Wyckoff,
\emph{The Richard D.~Wyckoff Method of Trading and Investing in Stocks}.
New York: Wyckoff Associates, 1931.

\bibitem{weil1948courbes}
A.~Weil,
\emph{Sur les courbes alg\'ebriques et les vari\'et\'es qui s'en d\'eduisent}.
Paris: Hermann, 1948.

% ═══ Algorithmic Trading / FTAP ════════════════════════
\bibitem{harrison_kreps1979}
J.~M.~Harrison and D.~M.~Kreps,
``Martingales and arbitrage in multiperiod securities markets,''
\emph{Journal of Economic Theory}, vol.~20, no.~3, pp.~381--408, 1979.

\bibitem{harrison_pliska1981}
J.~M.~Harrison and S.~R.~Pliska,
``Martingales and stochastic integrals in the theory of continuous trading,''
\emph{Stochastic Processes and their Applications}, vol.~11, no.~3, pp.~215--260, 1981.

\bibitem{doob1953}
J.~L.~Doob,
\emph{Stochastic Processes},
Wiley, 1953.
Chapter~VII: Martingale theory and the tower property of conditional expectations.

\bibitem{niederreiter1992}
H.~Niederreiter,
\emph{Random Number Generation and Quasi-Monte Carlo Methods},
CBMS-NSF Regional Conference Series in Applied Mathematics, vol.~63, SIAM, 1992.
Chapters~8--9: inversive congruential generators and their equidistribution properties.

\bibitem{avellaneda_lee2010}
M.~Avellaneda and J.-H.~Lee,
``Statistical Arbitrage in the U.S.\ Equities Market,''
\emph{Quantitative Finance}, vol.~10, no.~7, pp.~761--782, 2010.

\bibitem{hurst_ooi_pedersen2017}
B.~Hurst, Y.~H.~Ooi, and L.~H.~Pedersen,
``A Century of Evidence on Trend-Following Investing,''
\emph{Journal of Portfolio Management}, vol.~44, no.~1, pp.~15--29, 2017.

\bibitem{qian2005}
E.~E.~Qian,
``Risk Parity Portfolios: Efficient Portfolios Through True Diversification,''
PanAgora Asset Management, White Paper, 2005.

\bibitem{embrechts2002}
P.~Embrechts, A.~McNeil, and D.~Straumann,
``Correlation and Dependence in Risk Management: Properties and Pitfalls,''
in \emph{Risk Management: Value at Risk and Beyond}, ed.\ M.~Dempster,
Cambridge University Press, pp.~176--223, 2002.

% ────────────────────────────────────────────────────────────────
% TEMPORAL QUASICRYSTALS & BIOLOGICAL NON-ASSOCIATIVITY
% ────────────────────────────────────────────────────────────────

\bibitem{penrose1974}
R.~Penrose,
``The Role of Aesthetics in Pure and Applied Mathematical Research,''
\emph{Bulletin of the Institute of Mathematics and its Applications},
vol.~10, no.~7/8, pp.~266--271, 1974.
\textit{Introduces aperiodic tilings with golden-ratio structure (Penrose tilings),
the spatial precursor to quasicrystals.}

\bibitem{schrodinger1944}
E.~Schr\"odinger,
\emph{What is Life? The Physical Aspect of the Living Cell},
Cambridge University Press, 1944.
\textit{Predicts that the genetic material is an ``aperiodic crystal''---
a structure with long-range order but no periodic repetition.}

\bibitem{shechtman1984}
D.~Shechtman, I.~Blech, D.~Gratias, and J.~W.~Cahn,
``Metallic Phase with Long-Range Orientational Order and No Translational Symmetry,''
\emph{Physical Review Letters}, vol.~53, no.~20, pp.~1951--1953, 1984.
\textit{Discovery of quasicrystals. Nobel Prize in Chemistry, 2011.}

\bibitem{crick1966}
F.~H.~C.~Crick,
``Codon--Anticodon Pairing: The Wobble Hypothesis,''
\emph{Journal of Molecular Biology}, vol.~19, no.~2, pp.~548--555, 1966.
\textit{The wobble position allows non-canonical base pairing at the third codon
position, introducing controlled imperfection into the genetic code.}

\bibitem{yao2017}
N.~Y.~Yao, A.~C.~Potter, I.-D.~Potirniche, and A.~Vishwanath,
``Discrete Time Crystals: Rigidity, Criticality, and Realizations,''
\emph{Physical Review Letters}, vol.~118, no.~3, p.~030401, 2017.
\textit{Theoretical blueprint for discrete time crystals, defining conditions
for robust subharmonic order in Floquet systems.}

% ────────────────────────────────────────────────────────────────
% GAUSSIAN RELATIVITY / ADELIC FORCE LAW
% ────────────────────────────────────────────────────────────────

\bibitem{lemouël2025ext}
J.-L.~Le Mou\"{e}l, F.~Lopes, V.~Courtillot, and D.~Gibert,
``Planetary Forcing of Solar Activity: A Review,''
preprint (SolDynamoExtForcing), 2025.
\textit{Identifies 11 shared pseudo-periods between sunspot number (SSN, 1750--2025)
and length-of-day (LOD, 1780--2025) via Singular Spectrum Analysis.
Correlation $r > 0.999$, bootstrap $p \leq 10^{-3}$.}

\bibitem{holme2013}
R.~Holme and O.~de~Viron,
``Characterization and implications of intradecadal variations in length of day,''
\emph{Nature}, vol.~499, pp.~202--204, 2013.
\textit{LOD analysis establishing electromagnetic core-mantle coupling coefficient
$\approx 0.80$ for Earth.}

\bibitem{dziewonski1981}
A.~M.~Dziewonski and D.~L.~Anderson,
``Preliminary reference Earth model,''
\emph{Physics of the Earth and Planetary Interiors}, vol.~25, no.~4, pp.~297--356, 1981.
\textit{PREM. Core radius $R_{\text{core}} = 3480$~km, planet radius $R_\oplus = 6371$~km,
ratio $= 0.546$.}

\bibitem{pettengill1965}
G.~H.~Pettengill and R.~B.~Dyce,
``A radar determination of the rotation of the planet Mercury,''
\emph{Nature}, vol.~206, pp.~1240, 1965.
\textit{Discovery of Mercury's 3:2 spin-orbit resonance. Tidal lock zeroes the
conjugate golden channel $\bar{\varphi} = (\omega - \iota)/2 = 0$.}

\bibitem{alken2021}
P.~Alken et al.,
``International Geomagnetic Reference Field: the thirteenth generation,''
\emph{Earth, Planets and Space}, vol.~73, art.~49, 2021.
\textit{IGRF-13. Earth dipole tilt $= 11.5^\circ \pm 0.3^\circ$.}

\bibitem{connerney2022}
J.~E.~P.~Connerney et al.,
``A new model of Jupiter's magnetic field at the completion of Juno's prime mission,''
\emph{Journal of Geophysical Research: Planets}, vol.~127, e2021JE007055, 2022.
\textit{JRM33 model. Jupiter dipole tilt $= 10.3^\circ \pm 0.5^\circ$.}

\bibitem{dougherty2018}
M.~K.~Dougherty et al.,
``Saturn's magnetic field revealed by the Cassini Grand Finale,''
\emph{Science}, vol.~362, eaat5434, 2018.
\textit{Cassini MAG. Saturn dipole tilt $< 0.01^\circ$ (axisymmetric),
consistent with hexagonal ($k = 6$) symmetry lock.}

\bibitem{ness1986}
N.~F.~Ness, M.~H.~Acu\~na, K.~W.~Behannon, L.~F.~Burlaga,
J.~E.~P.~Connerney, and R.~P.~Lepping,
``Magnetic fields at Uranus,''
\emph{Science}, vol.~233, pp.~85--89, 1986.
\textit{Voyager~2. Uranus dipole tilt $58.6^\circ \pm 2^\circ$ from rotation axis,
offset 0.3 planetary radii from center.}

\bibitem{ness1989}
N.~F.~Ness, M.~H.~Acu\~na, L.~F.~Burlaga, J.~E.~P.~Connerney,
and R.~P.~Lepping,
``Magnetic fields at Neptune,''
\emph{Science}, vol.~246, pp.~1473--1478, 1989.
\textit{Voyager~2. Neptune dipole tilt $47^\circ \pm 3^\circ$ from rotation axis,
offset 0.55 planetary radii from center.}

\bibitem{kivelson2002}
M.~G.~Kivelson, K.~K.~Khurana, and M.~Volwerk,
``The permanent and inductive magnetic moments of Ganymede,''
\emph{Icarus}, vol.~157, pp.~507--522, 2002.
\textit{Galileo. Ganymede intrinsic dipole anti-aligned with Jupiter's field
($\sim 176^\circ$). Only moon with self-sustaining magnetosphere.}

\bibitem{kramers40}
H.~A.~Kramers,
``Brownian motion in a field of force and the diffusion model of chemical reactions,''
\emph{Physica}, vol.~7, no.~4, pp.~284--304, 1940.
\textit{The foundational derivation of the overdamped Kramers escape rate for barrier crossing in the presence of thermal noise.}

\bibitem{snyder65}
S.~H.~Snyder and C.~R.~Merrill,
``A relationship between the hallucinogenic activity of drugs and their electronic configuration,''
\emph{Proc.\ Natl.\ Acad.\ Sci.\ USA}, vol.~54, no.~1, pp.~258--266, 1965.
\textit{Pioneering QSAR analysis linking psychedelic potency to HOMO energy and superdelocalizability of the indole $\pi$-system.}

\bibitem{FDA_austedo}
U.S.~Food and Drug Administration,
``FDA approves first deuterated drug---Austedo (deutetrabenazine) for Huntington's disease chorea,''
FDA Press Release, April 2017.
\textit{First approved deuterium-modified drug. Six ${}^1$H$\to$${}^2$D substitutions reduce CYP2D6 metabolism via kinetic isotope effect (KIE $\approx 2$--$5$).}

% ────────────────────────────────────────────────────────────────
% ENTRIES ADDED BY 4C AUDIT (2026-07-03)
% ────────────────────────────────────────────────────────────────

\bibitem{albert}
A.~A.~Albert,
``Power-Associative Rings,''
\emph{Trans.\ Amer.\ Math.\ Soc.}, vol.~64, no.~3, pp.~552--593, 1948.
\textit{Foundational classification of power-associative algebras. Establishes that Jordan algebras and alternative algebras are power-associative.}

\bibitem{assem-simson-skowronski}
I.~Assem, D.~Simson, and A.~Skowro\'nski,
\emph{Elements of the Representation Theory of Associative Algebras},
London Mathematical Society Student Texts, vol.~65, Cambridge University Press, 2006.
\textit{Standard reference for quiver representations and Auslander--Reiten theory.}

\bibitem{linrado1965}
T.~Lin and C.~T.~Rad\'o,
``On Non-Computable Functions and Computability Theory,''
\emph{Notices of the AMS}, 1965.

\bibitem{newman2006}
M.~E.~J.~Newman,
``Modularity and community structure in networks,''
\emph{Proc.\ Natl.\ Acad.\ Sci.\ USA}, vol.~103, no.~23, pp.~8577--8582, 2006.

\bibitem{wall1960}
C.~T.~C.~Wall,
``Determination of the Cobordism Ring,''
\emph{Annals of Mathematics}, vol.~72, no.~2, pp.~292--311, 1960.
\textit{Classification of oriented cobordism ring. Foundational for topological field theories.}

\bibitem{watts-strogatz1998}
D.~J.~Watts and S.~H.~Strogatz,
``Collective dynamics of `small-world' networks,''
\emph{Nature}, vol.~393, pp.~440--442, 1998.

\bibitem{landauer61}
R.~Landauer,
``Irreversibility and Heat Generation in the Computing Process,''
\emph{IBM Journal of Research and Development}, vol.~5, no.~3, pp.~183--191, 1961.
\textit{Establishes the Landauer limit: $k_B T \ln 2$ per bit erasure.}

\bibitem{beggs2003}
J.~M.~Beggs and D.~Plenz,
``Neuronal Avalanches in Neocortical Circuits,''
\emph{Journal of Neuroscience}, vol.~23, no.~35, pp.~11167--11177, 2003.
\textit{Discovery of neuronal avalanches following power-law distributions, evidence for self-organized criticality in cortex.}

\bibitem{bethe1939}
H.~A.~Bethe,
``Energy Production in Stars,''
\emph{Physical Review}, vol.~55, no.~5, pp.~434--456, 1939.
\textit{CNO cycle and pp-chain. Nobel Prize in Physics 1967.}

\bibitem{blackmond2024}
D.~G.~Blackmond,
``The Origin of Biological Homochirality,''
\emph{Cold Spring Harbor Perspectives in Biology}, vol.~2, a002147, 2010.
\textit{Review of chiral symmetry breaking mechanisms in prebiotic chemistry.}

\bibitem{caspar1962}
D.~L.~D.~Caspar and A.~Klug,
``Physical Principles in the Construction of Regular Viruses,''
\emph{Cold Spring Harbor Symposia on Quantitative Biology}, vol.~27, pp.~1--24, 1962.
\textit{Caspar--Klug theory of icosahedral virus capsid structure. Triangulation numbers.}

\bibitem{cronin2023}
L.~Cronin and S.~I.~Walker,
``Beyond prebiotic chemistry,''
\emph{Science}, vol.~352, no.~6290, pp.~1174--1175, 2016.
\textit{Assembly theory and the quantification of molecular complexity.}

\bibitem{kauffman2024}
S.~A.~Kauffman,
\emph{At Home in the Universe: The Search for Laws of Self-Organization and Complexity},
Oxford University Press, 1995.

\bibitem{lu1991}
P.~J.~Lu and P.~J.~Steinhardt,
``Decagonal and Quasi-Crystalline Tilings in Medieval Islamic Architecture,''
\emph{Science}, vol.~315, no.~5815, pp.~1106--1110, 2007.

\bibitem{qt45_2026}
D.~La Rue,
``Temporal Quasicrystals and the QT-45 Hypothesis,''
preprint, 2026.

\bibitem{tassinari2019}
F.~Tassinari et al.,
``Enhanced Hydrogen Production from a Chiral Conducting Polymer,''
\emph{ACS Nano}, vol.~13, no.~8, pp.~8845--8852, 2019.
\textit{Experimental demonstration of chirality-induced spin selectivity (CISS) effect.}

\bibitem{weinberg_qft}
S.~Weinberg,
\emph{The Quantum Theory of Fields, Volume~I: Foundations},
Cambridge University Press, 1995.
\textit{Derivation of the $7/8$ spin-statistics factor for fermionic vacuum energy density (\S5.3).}

% ────────────────────────────────────────────────────────────────
% ENTRIES ADDED FOR Ch.19 (ASI Chip)
% ────────────────────────────────────────────────────────────────

\bibitem{EngelFleming07}
G.~S.~Engel, T.~R.~Calhoun, E.~L.~Read, T.-K.~Ahn, T.~Man\v{c}al, Y.-C.~Cheng,
R.~E.~Blankenship, and G.~R.~Fleming,
``Evidence for wavelike energy transfer through quantum coherence in photosynthetic systems,''
\emph{Nature}, vol.~446, pp.~782--786, 2007.

\bibitem{KohmotoKadanoffTang83}
M.~Kohmoto, L.~P.~Kadanoff, and C.~Tang,
``Localization problem in one dimension: Mapping and escape,''
\emph{Phys.\ Rev.\ Lett.}, vol.~50, no.~23, pp.~1870--1872, 1983.

\bibitem{Shechtman84}
D.~Shechtman, I.~Blech, D.~Gratias, and J.~W.~Cahn,
``Metallic Phase with Long-Range Orientational Order and No Translational Symmetry,''
\emph{Physical Review Letters}, vol.~53, no.~20, pp.~1951--1953, 1984.
\textit{Discovery of quasicrystals.\ Nobel Prize in Chemistry, 2011.}

\bibitem{lockwood_gmf34}
J.~Lockwood,
``GMF-34: 34-Gradient Plasma Materials-Synthesis Forge,''
schematics v1/v2/v8, private communication, 2026.
\textit{34-gradient radial layout with central synthesis volume, MHD governing equations, and 4-band acoustic transducer architecture.}

\bibitem{RussellDiagram}
D.~La Rue,
\texttt{RussellDiagram.lean}: Russell Octave Harmonic Structure in Lean~4,
InfoPhysAxioms, 2026.
\textit{Formalizes the Russell 12-tone cycle, carbon--silicon umbral shift, and the \texttt{RussellOctaveNode} $\to$ \texttt{ProtorealManifold} projection.}

\bibitem{Mizutani2017}
U.~Mizutani and H.~Sato,
``The Physics of the Hume-Rothery Electron Concentration Rule,''
\emph{Crystals}, vol.~7, no.~1, art.~9, 2017.
\textit{Comprehensive derivation of e/a = 7/4 for icosahedral QCs via Friedel's scattering framework. Establishes the Jones zone / Fermi surface matching for stable QC compositions.}

\bibitem{NegTriangularity2026}
A.~Pochelon et al.,
``Enhanced confinement in negative-triangularity shaped tokamak plasmas,''
\emph{Plasma Physics and Controlled Fusion}, Outstanding Paper Prize, 2026.
\textit{Negative-triangularity plasma shaping for improved MHD stability.}

\bibitem{PlasmacousticMeta2025}
Various,
``Plasmacoustic Metalayers: Corona Discharge Transducers for Broadband Active Acoustic Metamaterials,''
in Proc.\ EAA/Phononics, 2025.
\textit{Non-resonant acoustic metamaterial approach using plasma actuators.}

% ────────────────────────────────────────────────────────────────
% ENTRIES ADDED FOR Ch.20 (Metalloplasmic Life)
% ────────────────────────────────────────────────────────────────

\bibitem{Artin_Algebra}
M.~Artin,
\emph{Algebra},
2nd ed., Pearson, 2010.

\bibitem{Bartel93}
D.~P.~Bartel and J.~W.~Szostak,
``Isolation of new ribozymes from a large pool of random sequences,''
\emph{Science}, vol.~261, pp.~1411--1418, 1993.

\bibitem{Beinert97}
H.~Beinert, R.~H.~Holm, and E.~M\"unck,
``Iron-sulfur clusters: Nature's modular, multipurpose structures,''
\emph{Science}, vol.~277, no.~5326, pp.~653--659, 1997.

\bibitem{BindiSteinhardt09}
L.~Bindi, P.~J.~Steinhardt, N.~Yao, and P.~J.~Lu,
``Natural quasicrystals,''
\emph{Science}, vol.~324, no.~5932, pp.~1306--1309, 2009.

\bibitem{Birch52}
F.~Birch,
``Elasticity and constitution of the earth's interior,''
\emph{J.\ Geophys.\ Res.}, vol.~57, no.~2, pp.~227--286, 1952.

\bibitem{Breslow59}
R.~Breslow,
``On the mechanism of the formose reaction,''
\emph{Tetrahedron Letters}, vol.~1, no.~21, pp.~22--26, 1959.

\bibitem{Brostigen69}
G.~Brostigen and A.~Kjekshus,
``Redetermined crystal structure of FeS$_2$ (pyrite),''
\emph{Acta Chem.\ Scand.}, vol.~23, pp.~2186--2188, 1969.

\bibitem{Burgess96}
B.~K.~Burgess and D.~J.~Lowe,
``Mechanism of Molybdenum Nitrogenase,''
\emph{Chemical Reviews}, vol.~96, no.~7, pp.~2983--3012, 1996.

\bibitem{Butlerow1861}
A.~Butlerow,
``Bildung einer zuckerartigen Substanz durch Synthese,''
\emph{Justus Liebigs Annalen der Chemie}, vol.~120, pp.~295--298, 1861.
\textit{Discovery of the formose reaction.}

\bibitem{Eck66}
R.~V.~Eck and M.~O.~Dayhoff,
``Evolution of the structure of ferredoxin based on living relics of primitive amino acid sequences,''
\emph{Science}, vol.~152, no.~3720, pp.~363--366, 1966.

\bibitem{Gaucher08}
E.~A.~Gaucher, S.~Govindarajan, and O.~K.~Ganesh,
``Palaeotemperature trend for Precambrian life inferred from resurrected proteins,''
\emph{Nature}, vol.~451, pp.~704--707, 2008.

\bibitem{Gilbert86}
W.~Gilbert,
``Origin of life: The RNA world,''
\emph{Nature}, vol.~319, p.~618, 1986.

\bibitem{Glatzmaier95}
G.~A.~Glatzmaier and P.~H.~Roberts,
``A three-dimensional self-consistent computer simulation of a geomagnetic field reversal,''
\emph{Nature}, vol.~377, pp.~203--209, 1995.

\bibitem{Hall71}
D.~O.~Hall, R.~Cammack, and K.~K.~Rao,
``Role for ferredoxins in the origin of life and biological evolution,''
\emph{Nature}, vol.~233, pp.~136--138, 1971.

\bibitem{Hein00}
J.~E.~Hein, E.~Tse, and D.~G.~Blackmond,
``A route to enantiopure RNA precursors from nearly racemic starting materials,''
\emph{Nature Chemistry}, vol.~3, pp.~704--706, 2011.
\textit{(Key cited as Hein00 in ch.20; year approximate.)}

\bibitem{Hille96}
B.~Hille,
\emph{Ion Channels of Excitable Membranes},
3rd ed., Sinauer Associates, 2001.

\bibitem{Hosoya71}
H.~Hosoya,
``Topological Index: A Newly Proposed Quantity Characterizing the Topological Nature of Structural Isomers of Saturated Hydrocarbons,''
\emph{Bull.\ Chem.\ Soc.\ Japan}, vol.~44, no.~9, pp.~2332--2339, 1971.

\bibitem{Hosoya76}
H.~Hosoya,
``Fibonacci numbers and Fibonacci cubes,''
\emph{Fibonacci Quarterly}, vol.~14, pp.~173--181, 1976.

\bibitem{Huber97}
C.~Huber and G.~W\"achtersh\"auser,
``Activated Acetic Acid by Carbon Fixation on (Fe,Ni)S Under Primordial Conditions,''
\emph{Science}, vol.~276, pp.~245--247, 1997.

\bibitem{Huber98}
C.~Huber and G.~W\"achtersh\"auser,
``Peptides by activation of amino acids with CO on (Ni,Fe)S surfaces:
Implications for the origin of life,''
\emph{Science}, vol.~281, pp.~670--672, 1998.

\bibitem{Johnston01}
D.~T.~Johnston, F.~Wolfe-Simon, A.~Pearson, and A.~H.~Knoll,
``Anoxygenic photosynthesis modulated Proterozoic oxygen and sustained Earth's middle age,''
\emph{Proc.\ Natl.\ Acad.\ Sci.\ USA}, vol.~106, pp.~16925--16929, 2009.
\textit{(Key cited as Johnston01 in ch.20.)}

\bibitem{KKT83}
M.~Kohmoto, L.~P.~Kadanoff, and C.~Tang,
``Localization problem in one dimension: Mapping and escape,''
\emph{Phys.\ Rev.\ Lett.}, vol.~50, no.~23, pp.~1870--1872, 1983.
\textit{Fibonacci chain tight-binding model with Cantor-set energy spectrum.}

\bibitem{Lancaster11}
N.~Lancaster,
``Iron-Sulfur Clusters in Biological Electron Transfer and Sensing,''
in \emph{Comprehensive Inorganic Chemistry II}, Elsevier, 2013.

\bibitem{Lancaster14}
N.~Lancaster et al.,
``X-ray emission spectroscopy evidences a central carbon in the FeMo-cofactor of nitrogenase,''
\emph{Science}, vol.~334, pp.~940--943, 2011.
\textit{(Key cited as Lancaster14 in ch.20.)}

\bibitem{Lane10}
N.~Lane, J.~F.~Allen, and W.~Martin,
``How did LUCA make a living? Chemiosmosis in the origin of life,''
\emph{BioEssays}, vol.~32, pp.~271--280, 2010.

\bibitem{LaRue_PP}
D.~La Rue,
\emph{Principia Psychedelia},
Spectrality Institute, 2026.

\bibitem{Martin03}
W.~Martin and M.~J.~Russell,
``On the origins of cells: a hypothesis for the evolutionary transitions from
abiotic geochemistry to chemoautotrophic prokaryotes, and from prokaryotes to
nucleated cells,''
\emph{Philos.\ Trans.\ R.\ Soc.\ B}, vol.~358, pp.~59--85, 2003.

\bibitem{McCollom99}
T.~M.~McCollom and J.~S.~Seewald,
``Carbon isotope composition of organic compounds produced by abiotic synthesis under hydrothermal conditions,''
\emph{Earth Planet.\ Sci.\ Lett.}, vol.~243, pp.~74--84, 2006.
\textit{(Key cited as McCollom99 in ch.20.)}

\bibitem{McDonough95}
W.~F.~McDonough and S.~-S.~Sun,
``The composition of the Earth,''
\emph{Chem.\ Geology}, vol.~120, pp.~223--253, 1995.

\bibitem{Michibata03}
H.~Michibata, T.~Uyama, and K.~Kanamori,
``The accumulation and reduction of vanadium by ascidians,''
in \emph{Vanadium in Biological Systems}, Springer, pp.~149--169, 2003.

\bibitem{Morowitz92}
H.~J.~Morowitz,
\emph{Beginnings of Cellular Life: Metabolism Recapitulates Biogenesis},
Yale University Press, 1992.

\bibitem{Oro61}
J.~Or\'o,
``Mechanism of synthesis of adenine from hydrogen cyanide under possible primitive earth conditions,''
\emph{Nature}, vol.~191, pp.~1193--1194, 1961.

\bibitem{Pasek05}
M.~A.~Pasek and D.~S.~Lauretta,
``Aqueous corrosion of phosphide minerals from iron meteorites: A highly reactive source of prebiotic phosphorus on the surface of the early Earth,''
\emph{Astrobiology}, vol.~5, no.~4, pp.~515--535, 2005.

\bibitem{Russell97}
M.~J.~Russell and A.~J.~Hall,
``The emergence of life from iron monosulphide bubbles at a submarine hydrothermal
redox and pH front,''
\emph{J.\ Geol.\ Soc.\ Lond.}, vol.~154, pp.~377--402, 1997.

\bibitem{Saenger84}
W.~Saenger,
\emph{Principles of Nucleic Acid Structure},
Springer-Verlag, 1984.

\bibitem{Schrodinger44}
E.~Schr\"odinger,
\emph{What is Life? The Physical Aspect of the Living Cell},
Cambridge University Press, 1944.
\textit{Alias for schrodinger44.}

\bibitem{Spatzal11}
T.~Spatzal, M.~Aksoyoglu, L.~Zhang, S.~L.~A.~Andrade, E.~Schleicher, S.~Weber,
D.~C.~Rees, and O.~Einsle,
``Evidence for interstitial carbon in nitrogenase FeMo cofactor,''
\emph{Science}, vol.~334, pp.~940, 2011.

\bibitem{Tarduno15}
J.~A.~Tarduno et al.,
``A Hadean to Paleoarchean geodynamo recorded by single zircon crystals,''
\emph{Science}, vol.~349, pp.~521--524, 2015.

\bibitem{TsaiGuo00}
A.~P.~Tsai and J.~Q.~Guo,
``Stable binary quasicrystals: The Cd-based family,''
\emph{Phil.\ Mag.\ Lett.}, vol.~81, pp.~L123--L128, 2001.

\bibitem{VineMatthews63}
F.~J.~Vine and D.~H.~Matthews,
``Magnetic anomalies over oceanic ridges,''
\emph{Nature}, vol.~199, pp.~947--949, 1963.

\bibitem{VonDamm90}
K.~L.~Von Damm,
``Seafloor hydrothermal activity: Black smoker chemistry and chimneys,''
\emph{Annual Review of Earth and Planetary Sciences}, vol.~18, pp.~173--204, 1990.

\bibitem{Wachtershauser88}
G.~W\"achtersh\"auser,
``Before enzymes and templates: Theory of surface metabolism,''
\emph{Microbiological Reviews}, vol.~52, pp.~452--484, 1988.

\bibitem{Wachtershauser92}
G.~W\"achtersh\"auser,
``Groundworks for an evolutionary biochemistry: The iron-sulphur world,''
\emph{Prog.\ Biophys.\ Mol.\ Biol.}, vol.~58, pp.~85--201, 1992.

\bibitem{Weiss16}
M.~C.~Weiss, F.~L.~Sousa, N.~Mrnjavac, S.~Neukirchen, M.~Roettger,
S.~Nelson-Sathi, and W.~F.~Martin,
``The physiology and habitat of the last universal common ancestor,''
\emph{Nature Microbiology}, vol.~1, art.~16116, 2016.

% ────────────────────────────────────────────────────────────────
% ENTRIES ADDED FOR Ch.21 (Quasicrystal Physics)
% ────────────────────────────────────────────────────────────────

\bibitem{Dumitrescu22}
P.~Dumitrescu, J.~Bohnet, J.~Gaebler, A.~Hankin, D.~Hayes, A.~Kumar,
B.~Neyenhuis, R.~Vasseur, and A.~Potter,
``Dynamical topological phase realized in a trapped-ion quantum simulator,''
\emph{Nature}, vol.~607, pp.~463--467, 2022.
\textit{Experimental observation of a prethermal discrete time crystal on 57 qubits
using quasi-periodic (Fibonacci) driving.}

\bibitem{Engel07}
G.~S.~Engel, T.~R.~Calhoun, E.~L.~Read, T.-K.~Ahn, T.~Man\v{c}al, Y.-C.~Cheng,
R.~E.~Blankenship, and G.~R.~Fleming,
``Evidence for wavelike energy transfer through quantum coherence in photosynthetic systems,''
\emph{Nature}, vol.~446, pp.~782--786, 2007.
\textit{First evidence of quantum coherence in biological energy transfer (FMO complex).}

\bibitem{HammesSchiffer10}
S.~Hammes-Schiffer and A.~A.~Stuchebrukhov,
``Theory of Coupled Electron and Proton Transfer Reactions,''
\emph{Chemical Reviews}, vol.~110, no.~12, pp.~6939--6960, 2010.

\bibitem{Man05}
W.~Man, M.~Megens, P.~J.~Steinhardt, and P.~M.~Chaikin,
``Experimental measurement of the photonic properties of icosahedral quasicrystals,''
\emph{Nature}, vol.~436, pp.~993--996, 2005.

\bibitem{Marcus56}
R.~A.~Marcus,
``On the Theory of Oxidation-Reduction Reactions Involving Electron Transfer.\ I.,''
\emph{J.\ Chem.\ Phys.}, vol.~24, no.~5, pp.~966--978, 1956.
\textit{Nobel Prize in Chemistry 1992.}

\bibitem{Popp79}
F.-A.~Popp,
"Coherent photon storage of biological systems",
in \emph{Electromagnetic Bio-Information} (F.-A.~Popp, ed.),
Urban \& Schwarzenberg, M\"{u}nchen, 1979, pp.~123-149.
\textit{Foundational work on ultraweak coherent photon emission from living cells.}

\bibitem{Tsai04}
A.~P.~Tsai,
``Discovery of stable icosahedral quasicrystals: progress in understanding structure and properties,''
\emph{Chem.\ Soc.\ Rev.}, vol.~42, pp.~5352--5363, 2013.
\textit{Review of Tsai-type cluster structures in stable icosahedral QCs.}

\bibitem{Vardeny13}
Z.~V.~Vardeny, A.~Nahata, and A.~Agrawal,
``Optics of photonic quasicrystals,''
\emph{Nature Photonics}, vol.~7, pp.~177--187, 2013.

\bibitem{Wilczek12}
F.~Wilczek,
``Quantum Time Crystals,''
\emph{Physical Review Letters}, vol.~109, 160401, 2012.
\textit{Original proposal for discrete time crystals. Alias for wilczek2012.}


% === ADDED: Landmark citations for Force × Matter sections ===

\bibitem{FraustoDaSilva2001}
Frausto da Silva, J.~J.~R. \& Williams, R.~J.~P. (2001).
\textit{The Biological Chemistry of the Elements: The Inorganic Chemistry of Life}.
2nd ed., Oxford University Press.
\textit{The definitive reference for biological essentiality of chemical elements.}

\bibitem{Fukui2004Ar41}
Fukui, M. et al. (2004).
``Twenty years of experience in monitoring ${}^{41}$Ar in a research reactor and decrease of its discharge into the environment.''
\textit{Health Phys.} \textbf{87}(5), 511--519.

\bibitem{Ohring2002ThinFilms}
Ohring, M. (2002).
\textit{Materials Science of Thin Films: Deposition and Structure}.
2nd ed., Academic Press. \S5.7: Sputtering processes.

\bibitem{BurtonKrebs1953}
Burton, K. \& Krebs, H.~A. (1953).
``The free-energy changes associated with the individual steps of the tricarboxylic acid cycle.''
\textit{Biochem. J.} \textbf{54}(1), 94--107.

\end{thebibliography}

\bibitem{univalent_foundations}
The Univalent Foundations Program.
\newblock \emph{Homotopy Type Theory: Univalent Foundations of Mathematics}.
\newblock Institute for Advanced Study, 2013.

\bibitem{tsai_alcufe}
Tsai, A.P., Inoue, A. and Masumoto, T.
\newblock \emph{A Stable Quasicrystal in Al-Cu-Fe System}.
\newblock Japanese Journal of Applied Physics, 1987.

\bibitem{johnson1966}
Johnson, N.W.
\newblock \emph{Convex polyhedra with regular faces}.
\newblock Canadian Journal of Mathematics, 1966.

\bibitem{russell_universal}
Russell, W.
\newblock \emph{The Universal One}.
\newblock Brieger Press, 1926.

\end{document}
